\documentclass[11pt,letterpaper]{article}
\usepackage[T1]{fontenc}
\usepackage{amsmath,amsthm,mathtools}
\usepackage{newpxtext,newpxmath}
\usepackage[margin=1in,headheight=15pt]{geometry}
\usepackage{microtype}
\usepackage{xcolor,booktabs,array,enumitem}
\usepackage[most]{tcolorbox}
\usepackage{fancyhdr}
\usepackage{needspace}
\usepackage{hyperref}
\usepackage{bookmark}
\definecolor{Olive}{HTML}{48543A}
\definecolor{Pale}{HTML}{F3F4ED}
\definecolor{Ink}{HTML}{242924}
\hypersetup{colorlinks=true,linkcolor=Olive,citecolor=Olive,urlcolor=Olive,
 pdftitle={A Laurent-series case of Gonshor's birthday conjecture},
 pdfauthor={OpenAI ChatGPT},pdfsubject={Surreal numbers, ordinal natural products, exact birthday formulae}}
\usepackage{titlesec}
\titleformat{\section}{\Large\bfseries\color{Olive}}{\thesection}{0.7em}{}
\titleformat{\subsection}{\large\bfseries\color{Olive}}{\thesubsection}{0.7em}{}
\pagestyle{fancy}
\fancyhf{}
\fancyhead[L]{\small\scshape Surreal multiplication and ordinal length}
\fancyhead[R]{\small September 2026}
\fancyfoot[C]{\thepage}
\renewcommand{\headrulewidth}{0.3pt}
\setlength{\parskip}{0.35em}
\setlength{\parindent}{1.2em}
\setlist{nosep,leftmargin=1.6em}
\setcounter{tocdepth}{1}
\pretocmd{\section}{\Needspace{7\baselineskip}}{}{}
\pretocmd{\subsection}{\Needspace{5\baselineskip}}{}{}
\numberwithin{equation}{section}
\newtheorem{theorem}{Theorem}[section]
\newtheorem{lemma}[theorem]{Lemma}
\newtheorem{proposition}[theorem]{Proposition}
\newtheorem{corollary}[theorem]{Corollary}
\theoremstyle{definition}
\newtheorem{definition}[theorem]{Definition}
\newtheorem{example}[theorem]{Example}
\newtheorem{conjecture}[theorem]{Conjecture}
\theoremstyle{remark}
\newtheorem{remark}[theorem]{Remark}
\newcommand{\No}{\mathbf{No}}
\newcommand{\On}{\mathbf{On}}
\newcommand{\R}{\mathbb{R}}
\newcommand{\Z}{\mathbb{Z}}
\newcommand{\Q}{\mathbb{Q}}
\newcommand{\N}{\mathbb{N}}
\newcommand{\D}{\mathbb{Z}[1/2]}
\newcommand{\Lau}{\mathcal{L}}
\newcommand{\Fin}{\mathcal{F}}
\newcommand{\Pos}{\mathcal{P}}
\newcommand{\OrdPoly}{\mathcal{O}}
\newcommand{\bd}{\operatorname{b}}
\newcommand{\rd}{\operatorname{d}}
\newcommand{\supp}{\operatorname{supp}}
\newcommand{\tr}{\mathbin{\upharpoonright}}
\newcommand{\ns}{\mathbin{\oplus}}
\newcommand{\np}{\mathbin{\otimes}}
\newcommand{\bigns}{\mathop{\bigoplus}}
\newcommand{\co}[1]{\texttt{#1}}
\newtcolorbox{scopebox}{colback=Pale,colframe=Olive,boxrule=0.5pt,
 arc=1mm,left=9pt,right=9pt,top=7pt,bottom=7pt,breakable}

\begin{document}
\begin{center}
{\small\scshape Research article and reproducible proof audit}\par\vspace{0.8em}
{\LARGE\bfseries\color{Olive} A Laurent-series case of\par
Gonshor's birthday conjecture}\par\vspace{0.7em}
{\large Exact ordinal lengths, finite products, and infinite tails}\par\vspace{0.8em}
Prepared for Vladimir Reshetnikov\par
OpenAI ChatGPT\quad\textbullet\quad 20 September 2026
\end{center}
\vspace{0.6em}
\begin{abstract}
Gonshor's conjecture asks whether the birthday of a product of surreal numbers
is bounded by the natural product of their birthdays. We give a detailed proof
of this inequality on the canonically embedded Laurent-series field
$\R((\omega^{-1}))$. The argument first derives an exact birthday formula for
finite Laurent polynomials. The only information required from their bounded
part is the lowest exponent, its coefficient's birthday, and one dyadic versus
nondyadic boundary condition. A positive-polynomial/bounded-tail decomposition
then proves the product estimate for finite Laurent polynomials. A finite-jet
stabilization argument transfers this estimate to arbitrary Laurent series,
without assuming continuity of natural ordinal multiplication. We also prove
the conjecture for finite normal forms with ordinal exponents, determine the
birthday of every infinite Laurent series, and obtain exact birthdays of a
class of reciprocals. An exact-arithmetic program independently checks the
finite formula and explores fractional dyadic exponents beyond the proved
case.
\end{abstract}

\begin{scopebox}
\textbf{Status and scope.} This is a partial result, not a resolution of
Gonshor's conjecture for all surreal numbers. The proofs below use the standard
normal-form and sign-expansion theorems, with the subsequent arguments supplied
in detail. They have not been independently refereed or checked by a proof
assistant. The literature search located the documented conjecture and its
known monomial case, but did not establish priority for the Laurent-series
restriction proved here; no novelty claim is made.
\end{scopebox}

\newpage
\tableofcontents
\vspace{1.2em}

\section{The problem and the result}
\label{sec:problem}
A surreal number can be represented by a sequence of $+$ and $-$ signs indexed
by an ordinal. Its \emph{birthday}, written $\bd(x)$, is the length of this
sequence; in particular $\bd(0)=0$. For ordinals $\alpha,\beta$, let
$\alpha\ns\beta$ and $\alpha\np\beta$ denote their Hessenberg natural sum and
natural product. The problem selected for this investigation is the following.

\begin{conjecture}[Gonshor's birthday product conjecture]
\label{conj:gonshor}
For every $x,y\in\No$,
\begin{equation}\label{eq:conjecture}
 \bd(xy)\leq \bd(x)\np\bd(y).
\end{equation}
\end{conjecture}

Gonshor proposed the inequality in his 1986 book, p.~96
\cite{Gonshor1986}. Van den Dries and Ehrlich proved it when both factors are
single normal-form terms, $x=r\omega^a$ and $y=s\omega^b$, where $r,s$ are real
and $a,b$ are arbitrary surreal numbers \cite[Proposition~3.6]{vdDE2001}.
Their general estimate is the weaker bound
\begin{equation}\label{eq:knownweak}
 \bd(xy)\leq
 \omega\np\bd(x)\np\bd(x)\np\bd(y)\np\bd(y)
\end{equation}
\cite[Corollary~4.3]{vdDE2001}. In an August 2024 discussion, Ehrlich reported
that he knew of no improvement of the affirmative monomial case
\cite{Ehrlich2024}. Searches conducted for this article on 20 September 2026
did not locate a later resolution. That is evidence for selecting the problem,
not a proof of its current worldwide research status.

The principal result is a restriction of \eqref{eq:conjecture} to a familiar,
non-Archimedean field with unrestricted real coefficients.

\begin{theorem}[Laurent-series restriction]\label{thm:main}
Let
\begin{equation}\label{eq:field}
 \Lau=\R((\omega^{-1}))
 =\left\{\sum_{k\leq N}c_k\omega^k:
     N\in\Z,\ c_k\in\R\right\}\subseteq\No,
\end{equation}
where the sum is interpreted as a surreal normal form. Then
\[
 \bd(xy)\leq\bd(x)\np\bd(y)\qquad(x,y\in\Lau).
\]
\end{theorem}

The notation includes finite Laurent polynomials and series with infinitely
many nonzero terms. These are formal Hahn series, not sums taken in the
ordinary real topology. Integer exponents are bounded above, so every interval
of exponents bounded on both sides contains only finitely many possible terms.
This discreteness is crucial to the passage from finite to infinite products.

Two further results clarify the mechanism. For finite normal forms whose
exponents are ordinals, birthdays add by natural addition term by term
(Theorem~\ref{thm:ordinalpoly}). For an infinite Laurent series, however,
\begin{equation}\label{eq:infinite-preview}
 \bd(x)=\bd(x_{>0})+\omega^2,
\end{equation}
where the plus on the right is \emph{ordinary ordinal addition}
(Theorem~\ref{thm:infinite}). Confusing that plus with $\oplus$ would invalidate
a tempting direct extension of the finite proof.

\subsection{A guide to the proof}
Sections~\ref{sec:conventions}--\ref{sec:signs} establish the notation and the
sign-expansion input. Section~\ref{sec:ordinalpoly} handles ordinal-exponent
polynomials. Section~\ref{sec:finiteformula} computes finite Laurent birthdays
exactly. Section~\ref{sec:finiteproduct} proves the finite product estimate by
separating positive powers from the bounded tail. Section~\ref{sec:infinite}
passes to infinite series by preserving a fixed bound for every finite product
truncation. Sections~\ref{sec:examples}--\ref{sec:computation} give consequences,
failed generalizations, and reproducible checks.

\section{Three arithmetic conventions}
\label{sec:conventions}
The same objects occur both as surreal numbers and as ordinal lengths. We keep
the associated operations separate.

\subsection{Ordinary and natural ordinal arithmetic}
Unadorned $+$ between birthdays denotes \emph{ordinary ordinal addition}, the
operation corresponding to concatenation of well-ordered sequences. Thus
\[
 1+\omega=\omega,\qquad
 \omega+\omega^2=\omega^2,
 \qquad\omega^2+\omega>\omega^2.
\]
The symbols $\oplus$ and $\otimes$ always mean natural operations. If
\[
 \alpha=\sum_i\omega^{\xi_i}a_i,\qquad
 \beta=\sum_i\omega^{\xi_i}b_i
 \quad(\xi_0>\xi_1>\cdots)
\]
are Cantor normal forms written over the same finite list of exponents,
allowing zero coefficients, then
\[
 \alpha\ns\beta=\sum_i\omega^{\xi_i}(a_i+b_i).
\]
Natural multiplication is determined by distributivity and
\begin{equation}\label{eq:natmonomial}
 \omega^\xi\np\omega^\eta=\omega^{\xi\ns\eta}.
\end{equation}
Both operations are associative and commutative. Natural addition is strictly
increasing in either variable, and natural multiplication is strictly increasing
in either variable when the other is nonzero. In particular,
\[
 (\alpha\ns\beta)\np(\gamma\ns\delta)
 =\alpha\np\gamma\ns\alpha\np\delta
  \ns\beta\np\gamma\ns\beta\np\delta.
\]
These properties also follow from the Cantor-normal-form definitions, comparing
the largest exponent at which two expressions differ.

For ordinals below $\omega^\omega$, natural arithmetic is simply polynomial
arithmetic in a formal $\omega$, with nonnegative integer coefficients. For
example,
\[
 (\omega\cdot2+3)\np(\omega+4)
 =\omega^2\cdot2+\omega\cdot11+12.
\]
In an ordinal expression, $\omega^k\cdot n$ means ordinary right multiplication
by the finite integer $n$; it agrees with $\omega^k\otimes n$.

\subsection{Surreal arithmetic and normal forms}
In a surreal expression, such as $x+y$ or $\omega^a\omega^b$, the arithmetic
is that of the surreal field. The omega map obeys
\begin{equation}\label{eq:omega-field}
 \omega^a\omega^b=\omega^{a+b},
\end{equation}
where the exponent sum is surreal addition. When $a,b$ are ordinals this
surreal sum is their natural ordinal sum $a\oplus b$, not their ordinary sum.
When they are integers it is their usual integer sum.

We use the standard normal-form theorem: a surreal is represented uniquely by
\[
 x=\sum_{i<\lambda}r_i\omega^{a_i},
 \qquad a_0>a_1>\cdots,
 \qquad r_i\in\R\setminus\{0\},
\]
with reverse-well-ordered support. Addition and multiplication agree with Hahn
series operations \cite{Gonshor1986,BG2022}. All subsequent infinite series have
integer exponents bounded above; their nonzero supports have order type at most
$\omega$ when listed in decreasing order.

\subsection{A standard additive bound}
\begin{lemma}\label{lem:add}
For surreal $x,y$,
\[
 \bd(-x)=\bd(x),\qquad
 \bd(x+y)\leq\bd(x)\ns\bd(y).
\]
Consequently, $\bd(\sum_{j=1}^n x_j)\leq\bigns_{j=1}^n\bd(x_j)$.
\end{lemma}
\begin{proof}
Negation reverses every sign and preserves the length. For addition, use
canonical options, whose sign sequences are proper prefixes of the number's
sequence. Induct on $\bd(x)\oplus\bd(y)$. Every option in the defining cut for
$x+y$ has the form $x'+y$ or $x+y'$, where the replaced factor has smaller
birthday. By induction and strict monotonicity of natural addition, every such
option has birthday strictly below $\bd(x)\oplus\bd(y)$. A number defined by a
cut has birthday at most $\sup\{\bd(z)+1:z\text{ is an option}\}$.
Taking this supremum proves the claimed weak inequality. The finite version
follows by induction on $n$.
\end{proof}
This is the usual additive birthday theorem, not a new result; see
\cite[p.~96]{Gonshor1986}.

\section{Real coefficients and sign-expansion tools}
\label{sec:signs}
\subsection{Birthdays of real numbers}
Write $\rd(r)=\bd(r)$ for $r\in\R$. A real is \emph{dyadic} if it belongs to
$\D=\Z[1/2]$; zero is dyadic. Nonzero dyadics have finite birthdays, whereas
nondyadic reals have birthday $\omega$. More explicitly, for
$r\neq0$ dyadic with reduced denominator $2^h$,
\begin{equation}\label{eq:realbirth}
 \rd(r)=\lceil |r|\rceil+h.
\end{equation}
Thus $\rd(\pm1)=1$, $\rd(\pm\tfrac12)=2$, $\rd(\pm\tfrac34)=3$, and
$\rd(\tfrac13)=\omega$. Formula~\eqref{eq:realbirth} follows from the finite
binary sign expansion: for positive noninteger $r$, it starts with
$\lfloor r\rfloor+1$ plus signs and has $h$ further signs. Negation reverses
the signs. The formula also covers integers, for which $h=0$.

\begin{lemma}[Real scalar product bound]\label{lem:realproduct}
For all real $r,s$,
\[
 \rd(rs)\leq\rd(r)\np\rd(s).
\]
\end{lemma}
\begin{proof}
Zero factors are immediate. If at least one nonzero factor is nondyadic, the
right side is at least $\omega$, whereas $rs$ is real and has birthday at most
$\omega$. Suppose both factors are dyadic. Put
$A=\lceil|r|\rceil$, $B=\lceil|s|\rceil$, and let $2^h,2^k$ be their reduced
denominators. Then $A,B\geq1$, the reduced denominator exponent of $rs$ is at
most $h+k$, and $\lceil|rs|\rceil\leq AB$. Hence
\[
 \rd(rs)\leq AB+h+k
 \leq(A+h)(B+k)=\rd(r)\rd(s).
\]
Natural multiplication of finite ordinals is ordinary integer multiplication.
\end{proof}

\subsection{The sign-expansion input}
We use Gonshor's sign-expansion theorem in the formulation restated by Ehrlich
\cite[Propositions~3.7--3.8]{Ehrlich2011}; see also
\cite[Definition~2.6 and Theorem~2.7]{BG2022}. Here is the version needed for our
calculations.

If $a$ has a sign sequence and $p_\delta$ denotes the number of plus signs
strictly before its $\delta$th sign, the sequence for $\omega^a$ starts with a
plus. For each successive sign of $a$, append a block of that sign of length
$\omega^{p_\delta+1}$. For $r>0$, append the signs of $r$ after its first plus,
repeating each of them $\omega^p$ times, where $p$ is the total number of pluses
in $a$. Negative $r$ reverses the resulting signs.

For a normal form $\sum_i r_i\omega^{a_i}$, first \emph{reduce} each exponent's
sign sequence. A minus is deleted if its entire prefix, including that minus,
already occurs in an earlier exponent. There is a second rule: if the current
exponent starts with the immediately preceding exponent followed by a minus,
and the preceding coefficient is nondyadic, delete that additional minus.
Compute a term block from each reduced exponent by the preceding rule, then
concatenate the blocks. No plus signs are deleted.

\begin{lemma}[Prefixes and countable unions]\label{lem:prefix}
A normal-form truncation is a sign-sequence prefix, and therefore has birthday
at most that of the complete normal form. If a normal form has infinitely many
terms indexed by $\N$, then
\begin{equation}\label{eq:prefixsup}
 \bd\left(\sum_{i<\omega}r_i\omega^{a_i}\right)
 =\sup_{n<\omega}\bd\left(\sum_{i<n}r_i\omega^{a_i}\right).
\end{equation}
\end{lemma}
\begin{proof}
The blocks for an initial normal-form segment are unchanged by later terms:
both deletion rules look only at earlier exponents and coefficients. Thus its
sign sequence is the corresponding initial concatenation. At a limit of such
segments the complete sign sequence is their union, whose length is the
supremum of their lengths.
\end{proof}

\subsection{The negative-integer block calculation}
For an exponent consisting of $k$ minuses, with $k$ a nonnegative integer, let
$L(k,r)$ denote the birthday of the term block with nonzero real coefficient
$r$. Here the exponent is a \emph{reduced} exponent; it need not equal the
original exponent of a term.

\begin{lemma}\label{lem:block}
The exact block lengths are
\begin{equation}\label{eq:block}
L(k,r)=
\begin{cases}
 \rd(r),&k=0,\\
 \omega\cdot k+(\rd(r)-1),&k\geq1,\ r\in\D,\\
 \omega\cdot(k+1),&k\geq1,\ r\notin\D.
\end{cases}
\end{equation}
\end{lemma}
\begin{proof}
There are no plus signs in the exponent, so each exponent sign contributes
$\omega$ signs, and coefficient signs contribute one sign each. If $k\geq1$,
the initial plus is absorbed in ordinary ordinal addition:
$1+\omega\cdot k=\omega\cdot k$. A dyadic coefficient contributes precisely
$\rd(r)-1$ further signs. A nondyadic coefficient contributes $\omega$ further
signs. If $k=0$, only the real coefficient's sign sequence remains. Negating
$r$ changes no length.
\end{proof}

\section{An exact case with ordinal exponents}
\label{sec:ordinalpoly}
Consider the class ring
\[
 \OrdPoly=\left\{\sum_{i=1}^{m}r_i\omega^{\alpha_i}:
  m<\omega,\quad\alpha_i\in\On,\quad r_i\in\R\right\}.
\]
Zero coefficients are omitted and equal exponents are collected. It is a ring
because a product exponent is $\alpha_i\oplus\beta_j$, again an ordinal.
Normal forms supported on ordinals are necessarily finite: an infinite support
would yield an infinite strictly decreasing sequence of ordinals.

\begin{theorem}[Ordinal-exponent polynomials]\label{thm:ordinalpoly}
For $P=\sum_{i=1}^m r_i\omega^{\alpha_i}$ in decreasing normal form,
\begin{equation}\label{eq:ordinalbirth}
 \bd(P)=\bigns_{i=1}^m
       \bigl(\omega^{\alpha_i}\np\rd(r_i)\bigr).
\end{equation}
For all $P,Q\in\OrdPoly$,
\[
 \bd(PQ)\leq\bd(P)\np\bd(Q).
\]
\end{theorem}
\begin{proof}
An ordinal exponent has only plus signs, so neither deletion rule applies.
Since $\omega^\alpha$ is the corresponding ordinal, its birthday is
$\omega^\alpha$. The coefficient part of the sign expansion gives
\[
 \bd(r\omega^\alpha)=
 \begin{cases}
  \omega^\alpha\cdot\rd(r),&r\in\D\setminus\{0\},\\
  \omega^{\alpha+1},&r\notin\D.
 \end{cases}
\]
Both cases equal $\omega^\alpha\otimes\rd(r)$.

The birthday of the whole normal form is the ordinary concatenation sum of
these ordinal monomials. Their Cantor exponents are $\alpha_i$ or
$\alpha_i+1$. If $\alpha_i>\alpha_j$ then $\alpha_i\geq\alpha_j+1$, so those
Cantor exponents are nonincreasing in normal-form order. Consequently the
ordinary sum equals the natural sum, proving \eqref{eq:ordinalbirth}.

For two individual terms, Lemma~\ref{lem:realproduct} and
\eqref{eq:natmonomial} give
\begin{align*}
 \bd\bigl((r\omega^\alpha)(s\omega^\beta)\bigr)
 &=\omega^{\alpha\ns\beta}\np\rd(rs)\\
 &\leq(\omega^\alpha\np\rd(r))
        \np(\omega^\beta\np\rd(s)).
\end{align*}
Expand $PQ$ into its finitely many pairwise term products and apply
Lemma~\ref{lem:add}. Natural distributivity and
\eqref{eq:ordinalbirth} then give the result. Possible equal exponents or
coefficient cancellations can only improve the resulting upper bound; the
additive lemma does not require the expanded sum to be in normal form.
\end{proof}

\begin{example}
Let $P=3\omega^2+\tfrac13\omega-\tfrac12$. Its term birthdays are
$\omega^2\cdot3$, $\omega^2$, and $2$, respectively. Hence
\[
 \bd(P)=\omega^2\cdot4+2.
\]
The coefficient at exponent $1$ contributes at Cantor exponent $2$ because it
is nondyadic. This shift is why a degree estimate must distinguish dyadic and
nondyadic leading coefficients.
\end{example}

\section{The exact finite Laurent birthday formula}
\label{sec:finiteformula}
Let $\Fin=\R[\omega,\omega^{-1}]$. For $x\in\Fin$, write
\[
 x=P+u,
 \qquad P=x_{>0}=\sum_{k\geq1}c_k\omega^k,
 \qquad u=x_{\leq0}=\sum_{k\leq0}c_k\omega^k.
\]
Both sums are finite. The term \emph{bounded} means that the element lies
between two real numbers; for a Laurent polynomial this is equivalent to
having no positive exponent. Put
\begin{equation}\label{eq:H}
 H(P)=\bigns_{k\geq1}\bigl(\omega^k\np\rd(c_k)\bigr)=\bd(P),
\end{equation}
where zero coefficients contribute zero.

\begin{theorem}[Finite Laurent formula]\label{thm:finiteformula}
If $x$ has no negative-exponent term, then
\[
 \bd(x)=H(P)\ns\rd(c_0).
\]
Otherwise let $-m$, $m\geq1$, be its lowest nonzero exponent and put
$a=c_{-m}\neq0$. If $a\notin\D$, then
\begin{equation}\label{eq:finiteND}
 \bd(x)=H(P)+\omega\cdot(m+1).
\end{equation}
If $a\in\D$, define
\[
 \varepsilon=
 \begin{cases}
  1,&c_{-m+1}\notin\D,\\
  0,&c_{-m+1}\in\D,
 \end{cases}
\]
with an absent coefficient interpreted as zero. Then
\begin{equation}\label{eq:finiteD}
 \bd(x)=H(P)+\omega\cdot m+
             \bigl(\rd(a)-1+\varepsilon\bigr).
\end{equation}
All unadorned plus signs on these right sides are ordinary ordinal additions.
\end{theorem}

\begin{proof}
First compute the birthday of the bounded part $u$ alone. Its negative
exponents have sign sequences consisting entirely of minuses. Suppose the
current term is $c\omega^{-m}$ and the preceding nonzero term of the bounded
part has exponent $-p$, where $0\leq p<m$. Its sign sequence has $p$ minuses,
so the first deletion rule removes exactly those $p$ shared minuses. Put
$\chi=1$ if that preceding coefficient is nondyadic and $\chi=0$ otherwise.
The second rule removes one additional minus precisely when $\chi=1$.
Thus the current reduced exponent has
\begin{equation}\label{eq:gap}
 k=m-p-\chi
\end{equation}
minuses. This is nonnegative because $m-p\geq1$.

Before the first negative term, there may be a constant coefficient. In that
case the same description applies with $p=0$. If there is no constant term,
use an empty preceding part, $p=0$, $\chi=0$, and length zero.

We prove the claimed bounded-part formula inductively along its nonzero
negative terms. If the preceding coefficient is nondyadic, the preceding
length is $\omega\cdot(p+1)$. If it is dyadic, the preceding length is
$\omega\cdot p+t$ for a finite $t$; for an initial dyadic constant this simply
means its finite birthday. These statements are also valid for the empty
initial part with $p=t=0$.

Suppose first that $k\geq1$. Appending the block from
Lemma~\ref{lem:block} erases that finite $t$ and gives an $\omega$ coefficient
of
\[
 (p+\chi)+k=m
\]
when $c$ is dyadic, or $m+1$ when $c$ is nondyadic. In the dyadic case the
finite remainder is $\rd(c)-1$. In the nondyadic case it is zero.

The exceptional case is $k=0$. Equation~\eqref{eq:gap} then forces
$m=p+1$ and $\chi=1$: the previous coefficient is nondyadic at the immediately
adjacent exponent. The preceding length is $\omega\cdot m$, and the current
block has length $\rd(c)$. For dyadic $c$, the new length is therefore
$\omega\cdot m+\rd(c)$, one larger in its finite remainder than the usual
formula. For nondyadic $c$, it is $\omega\cdot(m+1)$. This proves the asserted
$\varepsilon$ correction and completes the induction.

It remains to restore $P$. Positive integer exponents have only plus signs.
They neither share a minus-containing prefix with a nonpositive exponent nor
trigger the second deletion rule across the positive/bounded boundary.
Consequently the sign sequence of $P+u$ is the sign sequence of $P$ followed
by exactly the same blocks as for $u$ alone. Its length is $H(P)+\bd(u)$.
If $u$ is real, Theorem~\ref{thm:ordinalpoly} gives the stated nonnegative
formula. Otherwise the calculation just made gives
\eqref{eq:finiteND} and \eqref{eq:finiteD}.
\end{proof}

\begin{corollary}[Exact finite splitting]\label{cor:split}
For finite Laurent polynomials $x=P+u$ as above,
\begin{equation}\label{eq:split}
 \bd(x)=\bd(P)\ns\bd(u).
\end{equation}
Every bounded finite Laurent polynomial has birthday below $\omega^2$. If it
has a negative-exponent term, its birthday is at least $\omega$.
\end{corollary}
\begin{proof}
The finite formula gives $\bd(u)=\omega\cdot M+\tau$ with $M,\tau$ finite,
or a real birthday at most $\omega$. On the other hand, every nonzero Cantor
term of $H(P)$ has exponent at least $1$. Hence concatenating $\bd(u)$ after
$H(P)$ discards no nonzero Cantor term of $H(P)$. Ordinary addition and natural
addition agree in this specific configuration. The other assertions follow
immediately from the formula.
\end{proof}

\subsection{Why the boundary correction matters}
The exact value of the coefficient immediately above the lowest exponent is
irrelevant, except for whether it is nondyadic. For example,
\begin{align*}
 \bd(1+\omega^{-1})&=\omega,&
 \bd(\tfrac13+\omega^{-1})&=\omega+1,\\
 \bd(\tfrac13+\omega^{-2})&=\omega\cdot2,&
 \bd(\tfrac13\omega^{-1}+\tfrac12\omega^{-2})
 &=\omega\cdot2+2.
\end{align*}
In the last example, the second negative exponent is reduced to the empty
sequence, so its coefficient contributes two signs rather than one. A rule
that always appends $\rd(a)-1$ signs would miss this case. A gap of two or more
integer exponents prevents this finite correction from surviving, although
the nondyadic deletion still takes place within the intermediate calculation.

\section{The finite product theorem}
\label{sec:finiteproduct}
The exact splitting formula reduces the problem to three products: positive
parts with positive parts, positive parts with bounded parts, and bounded
parts with bounded parts. The first was covered by
Theorem~\ref{thm:ordinalpoly}.

\subsection{Bounded factors}
\begin{lemma}\label{lem:bounded}
For bounded finite Laurent polynomials $u,v$,
\[
 \bd(uv)\leq\bd(u)\np\bd(v).
\]
If $u$ has a negative-exponent term and $v\notin\{0,1,-1\}$, this inequality
is strict.
\end{lemma}
\begin{proof}
Zero factors cause no difficulty. If both factors are real, apply
Lemma~\ref{lem:realproduct}. If both have negative-exponent terms, then
\[
 \bd(uv)<\omega^2\leq\bd(u)\np\bd(v),
\]
since their product is still a bounded finite Laurent polynomial and both
factor birthdays are at least $\omega$.

It remains to multiply a bounded finite Laurent polynomial $u$ with a genuine
negative term by a nonzero real $r$. Write
\[
 \bd(u)=\omega\cdot M+\tau,\qquad M\geq1,\quad M,\tau<\omega.
\]
If $r$ is nondyadic, the proposed bound is at least $\omega^2$, while
$\bd(ru)<\omega^2$. If $r$ is dyadic, put $d=\rd(r)$. Multiplication by a
nonzero real leaves the lowest exponent of $u$ unchanged. A dyadic lowest
coefficient stays dyadic; a nondyadic one can stay nondyadic or become dyadic.
By the exact formula, the leading $\omega$ coefficient of $\bd(ru)$ is thus at
most $M$. For $d\geq2$, the proposed bound has leading coefficient $dM>M$:
\[
 \bd(u)\np d=\omega\cdot(dM)+d\tau.
\]
This proves strictness. The only nonzero reals with $d=1$ are $r=\pm1$;
then $\bd(ru)=\bd(u)$.
\end{proof}

\subsection{A coarse upper bound with one sharp edge case}
\begin{lemma}\label{lem:degree}
Let $z$ be a finite Laurent polynomial with all exponents at most $n$, where
$n\geq1$ is an integer. Then
\[
 \bd(z)<\omega^{n+1}\cdot2.
\]
More precisely, if the coefficient at exponent $n$ is dyadic, including zero,
then $\bd(z)<\omega^{n+1}$. If it is nondyadic, then
\[
 \bd(z)=\omega^{n+1}+\rho
 \qquad\text{for some }\rho<\omega^{n+1}.
\]
\end{lemma}
\begin{proof}
By Theorem~\ref{thm:finiteformula}, the positive part contributes at most one
term at Cantor exponent $n+1$, arising only from a nondyadic coefficient at
exponent $n$. That term has coefficient one. All remaining positive-power
contributions have Cantor exponent at most $n$, and the bounded contribution
is below $\omega^2$. Since $n\geq1$, their finite natural sum is below
$\omega^{n+1}$.
\end{proof}

\begin{lemma}[Positive monomial times bounded part]\label{lem:crossmono}
Let $n\geq1$ be an integer, $r\neq0$ real, and $u$ a bounded finite Laurent
polynomial. Then
\begin{equation}\label{eq:crossmono}
 \bd(r\omega^n u)\leq
       \bd(r\omega^n)\np\bd(u).
\end{equation}
If $u$ has a negative-exponent term, the inequality is strict.
\end{lemma}
\begin{proof}
If $u$ is real, both factors belong to $\OrdPoly$, and
Theorem~\ref{thm:ordinalpoly} applies. Otherwise write
$\bd(u)=\omega\cdot M+\tau$, with $M\geq1$ and $M,\tau$ finite.
The product $r\omega^n u$ has highest exponent at most $n$.

If $r$ is nondyadic, then $\bd(r\omega^n)=\omega^{n+1}$, and the right side
of \eqref{eq:crossmono} has leading Cantor exponent $n+2$. Lemma~\ref{lem:degree}
bounds the left side by an ordinal with leading exponent at most $n+1$, so
strictness follows.

Suppose that $r$ is dyadic and write $d=\rd(r)$. The proposed bound is
\begin{equation}\label{eq:crossbudget}
 \omega^{n+1}\cdot(dM)+\omega^n\cdot(d\tau).
\end{equation}
If $dM\geq2$, Lemma~\ref{lem:degree} makes the left side strictly smaller
already. Only the case $d=M=1$ remains.

Here $r=\pm1$. The finite formula forces the lowest exponent of $u$ to be
$-1$, with a nonzero dyadic coefficient $a$. Thus
\[
 u=c+a\omega^{-1},\qquad c\in\R,
 \qquad a\in\D\setminus\{0\}.
\]
If $c$ is dyadic, the top coefficient of $r\omega^n u$ is dyadic, so the
left side is below $\omega^{n+1}$ and the result follows. If $c$ is
nondyadic, the boundary correction is indispensable:
\[
 \bd(u)=\omega+\rd(a).
\]
The ordinal-exponent formula gives
\begin{align*}
 \bd(r\omega^n u)
   &=\omega^{n+1}+\omega^{n-1}\cdot\rd(a),\\
 \bd(r\omega^n)\np\bd(u)
   &=\omega^{n+1}+\omega^n\cdot\rd(a).
\end{align*}
Since $n\geq1$ and $\rd(a)>0$, the first is strictly smaller. This includes
$n=1$, when the lower term on the first line is finite.
\end{proof}

\begin{corollary}\label{cor:crosspoly}
For a positive-power polynomial $P$ and a bounded finite Laurent polynomial
$u$,
\[
 \bd(Pu)\leq\bd(P)\np\bd(u).
\]
If $P\neq0$ and $u$ has a negative-exponent term, the inequality is strict.
\end{corollary}
\begin{proof}
Expand $P$ into its finitely many nonzero terms. Apply
Lemma~\ref{lem:add} to the products, Lemma~\ref{lem:crossmono} to each term,
and Theorem~\ref{thm:ordinalpoly} to recombine the birthdays of $P$ by natural
addition. For strictness, at least one strict term inequality occurs, and
natural addition is strictly increasing.
\end{proof}

\subsection{Putting the four products together}
\begin{theorem}[Finite Laurent product theorem]\label{thm:finiteproduct}
For all $x,y\in\Fin$,
\[
 \bd(xy)\leq\bd(x)\np\bd(y).
\]
\end{theorem}
\begin{proof}
Write $x=P+u$ and $y=Q+v$ with positive-power polynomials $P,Q$ and bounded
parts $u,v$. Expand
\[
 xy=PQ+Pv+Qu+uv.
\]
By the additive birthday lemma and the three product estimates,
\begin{align*}
 \bd(xy)
 &\leq\bd(PQ)\ns\bd(Pv)\ns\bd(Qu)\ns\bd(uv)\\
 &\leq\bigl(\bd(P)\np\bd(Q)\bigr)
     \ns\bigl(\bd(P)\np\bd(v)\bigr)\\
 &\hspace{1.2em}\ns\bigl(\bd(u)\np\bd(Q)\bigr)
     \ns\bigl(\bd(u)\np\bd(v)\bigr)\\
 &=\bigl(\bd(P)\ns\bd(u)\bigr)
     \np\bigl(\bd(Q)\ns\bd(v)\bigr)\\
 &=\bd(x)\np\bd(y),
\end{align*}
where the last equality is the exact splitting formula
\eqref{eq:split}. This proof does not assume that any of the four product
summands has a disjoint support from the others.
\end{proof}

\begin{corollary}[Equality with a finite negative tail]\label{cor:strict}
If $x\in\Fin$ has a negative-exponent term, then for $y\in\Fin$,
\[
 \bd(xy)=\bd(x)\np\bd(y)
 \quad\Longleftrightarrow\quad y\in\{0,1,-1\}.
\]
\end{corollary}
\begin{proof}
The three displayed values give equality directly. Otherwise, in the preceding
proof $u$ has a negative-exponent term. If $Q\neq0$, the $Qu$ estimate is
strict by Corollary~\ref{cor:crosspoly}. If $Q=0$, then $y=v$ is bounded and
is not $0$ or $\pm1$, so the $uv$ estimate is strict by
Lemma~\ref{lem:bounded}. Strict monotonicity of the finite natural sum carries
that strictness to the complete bound.
\end{proof}

\section{Infinite Laurent series and finite jets}
\label{sec:infinite}
\subsection{The field and its product law}
With $t=\omega^{-1}$, the field in \eqref{eq:field} is the usual formal Laurent
series field $\R((t))$ embedded by normal forms. For completeness, its field
property can be seen directly. Every nonzero element has the shape
\[
 c\,t^m(1+h),\qquad c\in\R\setminus\{0\},\quad
 m\in\Z,\quad h\in t\R[[t]].
\]
Its inverse is the formal series
\[
 c^{-1}t^{-m}\sum_{j\geq0}(-h)^j.
\]
For any fixed power of $t$, only finitely many terms contribute, since every
term of $h$ has exponent at least one. The coefficientwise geometric-series
identity proves that this is an inverse.

In omega-exponent notation, if $x$ and $y$ have highest exponents $N_x$ and
$N_y$, the coefficient at exponent $k$ in their product is
\[
 [\omega^k](xy)=\sum_{i+j=k}c_i d_j.
\]
This sum is finite: the possible $i$ satisfy $k-N_y\leq i\leq N_x$.
The argument uses formal coefficients only, not analytic convergence.

\subsection{Exact birthdays of infinite tails}
\begin{theorem}[Infinite Laurent birthday formula]\label{thm:infinite}
Suppose $x\in\Lau$ has infinitely many nonzero terms. Let
$P=x_{>0}$ and $H=\bd(P)$. Then
\begin{equation}\label{eq:infinite}
 \boxed{\bd(x)=H+\omega^2,}
\end{equation}
with ordinary ordinal addition. In particular, every bounded infinite Laurent
series has birthday exactly $\omega^2$.
\end{theorem}
\begin{proof}
The positive part is finite. List the nonzero negative exponents as
$-m_1>-m_2>\cdots$. The positive integers $m_j$ tend to infinity, since the
support is infinite and integer-valued. Let $x_j$ be the finite normal-form
truncation through exponent $-m_j$.
Theorem~\ref{thm:finiteformula} gives
\[
 \bd(x_j)=H+\omega\cdot M_j+\tau_j,
\]
where $M_j$ equals $m_j$ or $m_j+1$, and $\tau_j$ is finite. Consequently these
lengths are cofinal in $H+\omega^2$. The sign sequences are nested prefixes by
Lemma~\ref{lem:prefix}, and their union is the sign sequence of $x$.
Taking their supremum proves \eqref{eq:infinite}.
\end{proof}

\begin{remark}[A real loss of a Cantor term]\label{rem:absorption}
Let $x=\omega+\sum_{n\geq1}\omega^{-n}$. Then
\[
 \bd(x)=\omega+\omega^2=\omega^2,
 \qquad
 \bd(\omega)\ns\bd\left(\sum_{n\geq1}\omega^{-n}\right)
 =\omega^2+\omega.
\]
Thus the finite splitting equality \eqref{eq:split} does not extend unchanged
to infinite tails. The difference is not a notation preference: it is the
ordinary-ordinal absorption of a genuine $\omega$ term.
\end{remark}

\subsection{Stabilization of a finite product truncation}
For an integer $K$, define
\[
 x_{\geq K}=\sum_{k\geq K}c_k\omega^k.
\]
This is a finite normal-form truncation for every $x\in\Lau$.

\begin{lemma}[Finite-jet stabilization]\label{lem:jets}
Let $x,y\in\Lau\setminus\{0\}$ have highest exponents $N_x,N_y$.
Fix $K\in\Z$ and choose any integer $T$ such that
\begin{equation}\label{eq:T}
 T\leq K-N_y,\qquad T\leq K-N_x.
\end{equation}
Then
\begin{equation}\label{eq:jet}
 (x_{\geq T}y_{\geq T})_{\geq K}=(xy)_{\geq K}.
\end{equation}
\end{lemma}
\begin{proof}
A discarded term of $x$ has exponent $i<T$. Multiplied by any term of $y$,
it has exponent $i+j<T+N_y\leq K$. Likewise, a discarded term of $y$ cannot
contribute to exponent $K$ or above. Therefore the complete products have the
same coefficients at every exponent at least $K$. These coefficients are
finite sums, so there is no exchange of infinite summation operations hidden
in the argument.
\end{proof}

\begin{proof}[Proof of Theorem~\ref{thm:main}]
Zero factors are immediate. Fix nonzero $x,y\in\Lau$ and put
\[
 B=\bd(x)\np\bd(y).
\]
For any integer cutoff $K$, choose $T$ as in Lemma~\ref{lem:jets}. Since normal
truncations are sign prefixes, and since $x_{\geq T},y_{\geq T}$ are finite
Laurent polynomials, we have
\begin{align*}
 \bd\bigl((xy)_{\geq K}\bigr)
 &=\bd\bigl((x_{\geq T}y_{\geq T})_{\geq K}\bigr)\\
 &\leq\bd(x_{\geq T}y_{\geq T})\\
 &\leq\bd(x_{\geq T})\np\bd(y_{\geq T})\\
 &\leq\bd(x)\np\bd(y)=B.
\end{align*}
The middle inequality is Theorem~\ref{thm:finiteproduct}, and the last one is
monotonicity of natural multiplication.

If $xy$ has finite support, choose $K$ below its lowest exponent, and this is
the desired bound. If it has infinite support, the finite truncations exhaust
its sign sequence by Lemma~\ref{lem:prefix}. Their birthdays are all bounded
by the \emph{same} ordinal $B$, so their supremum is bounded by $B$ as well.
Thus $\bd(xy)\leq B$ in either case.
\end{proof}

The proof deliberately does not assert a general continuity identity such as
$\sup_n(\alpha_n\otimes\beta_n)
=(\sup_n\alpha_n)\otimes(\sup_n\beta_n)$.
It only keeps one fixed upper bound throughout. Even natural addition need
not preserve a supremum in a variable: $\sup_{n<\omega}(n\oplus1)=\omega$,
whereas $(\sup_{n<\omega}n)\oplus1=\omega+1$.

\Needspace{11\baselineskip}
\section{Examples and exact consequences}
\label{sec:examples}
\subsection{A table of birthdays}
The following values follow directly from the finite and infinite formulas.
The sums displayed here are surreal sums; every expression in the right column
is an ordinal in Cantor normal form.

\begin{center}
\renewcommand{\arraystretch}{1.32}
\begin{tabular}{@{}p{0.65\textwidth}l@{}}
\toprule
Surreal number & Birthday\\
\midrule
$1+\omega^{-1}$ & $\omega$\\
$\tfrac13+\omega^{-1}$ & $\omega+1$\\
$\omega+\omega^{-1}$ & $\omega\cdot2$\\
$\omega^{-m}$, $m\geq1$ & $\omega\cdot m$\\
$\tfrac12\omega^{-m}$, $m\geq1$ & $\omega\cdot m+1$\\
$\tfrac13\omega^{-m}$, $m\geq1$ & $\omega\cdot(m+1)$\\
$\tfrac13\omega^{-1}+\tfrac12\omega^{-2}$ & $\omega\cdot2+2$\\
$\sum_{n\geq1}\omega^{-n}$ & $\omega^2$\\
$\omega+\sum_{n\geq1}\omega^{-n}$ & $\omega^2$\\
$\omega^3+\sum_{n\geq1}\omega^{-n}$ & $\omega^3+\omega^2$\\
$\tfrac13\omega^2+\sum_{n\geq1}\omega^{-n}$ & $\omega^3+\omega^2$\\
$\sum_{n\geq0}\omega^{-n}/n!$ & $\omega^2$\\
\bottomrule
\end{tabular}
\end{center}

\subsection{How loose can the natural-product bound be?}
For positive integers $m,n$,
\[
 \bd(\omega^{-m}\omega^{-n})=\omega\cdot(m+n),
 \qquad
 \bd(\omega^{-m})\np\bd(\omega^{-n})=\omega^2\cdot mn.
\]
The conjectured bound is therefore already very non-sharp for elementary
infinitesimals. Cancellation can make the discrepancy larger still. If
\[
 x=\sum_{n\geq0}\omega^{-n},\qquad y=1-\omega^{-1},
\]
then $xy=1$, but
\[
 \bd(x)=\omega^2,\qquad\bd(y)=\omega,
 \qquad\bd(x)\np\bd(y)=\omega^3.
\]
The theorem concerns a uniform complexity bound, not an estimate expected to
be nearly an equality for typical non-Archimedean factors.

\subsection{Reciprocals of finite Laurent polynomials}
\begin{corollary}\label{cor:reciprocal}
Let $f\in\Fin\setminus\{0\}$ have at least two nonzero terms and highest
exponent $N\geq0$. Then
\begin{equation}\label{eq:reciprocal}
 \bd(f^{-1})=\omega^2.
\end{equation}
\end{corollary}
\begin{proof}
The only units in the finite Laurent polynomial ring over a field are the
nonzero monomials. Indeed, for a nonzero finite Laurent polynomial define its
width as highest exponent minus lowest exponent. The extreme coefficients of
a product are nonzero products of extreme coefficients, so widths add under
multiplication. A product equal to $1$ has width zero, forcing both factors to
have width zero.

Thus $f^{-1}$ is not finite; it has infinitely many terms in $\Lau$. Its
highest exponent is $-N\leq0$, so it is bounded. Apply
Theorem~\ref{thm:infinite}.
\end{proof}

For example,
\[
 (\omega+1)^{-1}
 =\omega^{-1}-\omega^{-2}+\omega^{-3}-\cdots
=\sum_{n\geq1}(-1)^{n-1}\omega^{-n}
\]
where the displayed identity is coefficientwise formal multiplication; its
birthday is $\omega^2$. Likewise, $(1+\omega^{-1})^{-1}$ and
$(\omega^2+\omega)^{-1}$ have birthday $\omega^2$.
The assumption $N\geq0$ cannot be silently omitted: a reciprocal with positive
leading exponent can have a positive polynomial part that survives in
$H+\omega^2$.

\subsection{The complete birthday spectrum of this field}
\begin{corollary}\label{cor:spectrum}
The set of birthdays realized in $\Lau$ is exactly
\[
 \{\alpha:\alpha<\omega^\omega\}.
\]
\end{corollary}
\begin{proof}
The finite and infinite formulas place every birthday below $\omega^\omega$.
Conversely, every ordinal $\alpha<\omega^\omega$ has a finite Cantor normal
form with finite exponents and nonnegative integer coefficients. It is itself
an element of $\Lau$, and $\bd(\alpha)=\alpha$.
\end{proof}

This does not identify $\Lau$ with the entire set of surreals of birthday
below $\omega^\omega$. For instance, $\omega^{1/2}$ has birthday $\omega^2$
but is not in $\Lau$, because its unique normal form has a noninteger exponent.
In particular $\Lau$ is not real closed: its positive element $\omega$ has no
square root in $\Lau$.

\section{What the proof does not extend to automatically}
\label{sec:limits}
\subsection{The failed termwise shortcut}
The known monomial case might suggest expanding every product into term
products and summing their birthday bounds. This gives
\[
 \bd(xy)\leq
 \left(\bigns_i\bd(r_i\omega^{a_i})\right)
 \np
 \left(\bigns_j\bd(s_j\omega^{b_j})\right)
\]
for finite normal forms, but does not in general give the desired bound in
terms of $\bd(x)$ and $\bd(y)$. One would need equality between a number's
birthday and the natural sum of its term birthdays, which can fail.

The simplest Laurent example is
\[
 \bd(1+\omega^{-1})=\omega
 <\omega+1=\bd(1)\ns\bd(\omega^{-1}).
\]
Even with positive fractional exponents,
\[
 \bd(\omega+\omega^{1/2})=\omega^2,
 \qquad
 \bd(\omega)\ns\bd(\omega^{1/2})=\omega^2+\omega.
\]
Here $1/2$ has sign sequence $+-$, so $\omega^{1/2}$ has a first block of
$\omega$ pluses followed by $\omega^2$ minuses. Its birthday is $\omega^2$.
The leading $\omega$ term of the sum is absorbed in the concatenated birthday.

The positive/bounded decomposition in the finite Laurent proof was chosen
precisely to retain exact equality at the recombination step. It is not a
claim that every normal-form decomposition has that property.

\subsection{Fractional dyadic exponents}
A natural next restriction is
\[
 \R[\omega^q:q\in\D],
\]
where each expression has finite support. The general sign-expansion algorithm
still works with finite exponent sign sequences, so exact computations are
feasible. What is missing is an analogue of the finite Laurent formula that
retains the sharp birthday budget under arbitrary fractional-exponent
products.

The obstacle is structural. Positive integer exponents have no minuses, while
negative integer exponents lie on a single all-minus branch. This made the
reduction rule depend only on the previous negative exponent and one
coefficient flag. A general dyadic exponent can have an arbitrary finite
mixture of signs. An earlier, nonadjacent exponent can then delete a minus in
a later exponent, and multiple sign-block scales can interact. The last
exponent and its immediate predecessor no longer summarize all relevant
information.

The accompanying search found no counterexample in 100,000 random pairs from
a specified dyadic grid, but this is not a proof of that larger restriction.
A useful next technical target is a \emph{prefix-sensitive} bound on the reduced
term blocks: it must account for the finite tree of shared exponent prefixes,
rather than merely adding independent term lengths.

\subsection{Infinite supports beyond the integer lattice}
The finite-jet argument used that every cutoff $x_{\geq K}$ is finite. This is
not true for arbitrary reverse-well-ordered supports in a dense exponent group.
For example, the descending sequence
\[
 -\frac12,-\frac34,-\frac78,-\frac{15}{16},\ldots
\]
is a legitimate dyadic support, but all its infinitely many terms lie above
$-1$. Thus the same numerical cutoff no longer produces a finite polynomial.
This example does not disprove a more general approximation theorem; it marks
the specific step that has not been justified outside the discrete Laurent
setting.

\subsection{No general solution has been obtained}
The full conjecture permits arbitrary surreal exponents, arbitrary set-sized
reverse-well-ordered supports, and products for which both sign reductions and
support order types are much more complicated. Neither the polynomial theorem
with ordinal exponents nor the Laurent-field theorem covers those cases.
The computational search supplies evidence, not the missing argument.

\section{Exact computation and reproducibility}
\label{sec:computation}
The archive includes \co{verify.py}, a standard-library-only Python program
compatible with Python 3.9 or later, and the recorded results in
\co{verification\_results.json}. Mathematical arithmetic is exact: coefficients
and dyadic exponents use \co{fractions.Fraction}; ordinal operations use integer
coefficient tuples. No floating-point approximation enters a birthday comparison.

\subsection{Two independent routes to finite Laurent birthdays}
The function \co{sign\_birthday} implements the general finite-exponent
sign-expansion theorem. For each exponent it records its finite sign sequence,
removes already-shared minus prefixes, applies the separate nondyadic
predecessor deletion, and concatenates ordinal block lengths.

The function \co{laurent\_formula} instead implements
Theorem~\ref{thm:finiteformula}, using only positive-power term birthdays, the
lowest exponent and coefficient, and the single boundary flag. It does not call
the sign-expansion engine. Agreement between the two implementations is a
check on the algebraic simplification and on edge-case handling, not an
independent proof of Gonshor's sign-expansion theorem.

For instance, the tuple $(2,4,0,1)$ represents
$\omega^3+\omega\cdot4+2$. Natural addition is componentwise addition, natural
multiplication is convolution, and ordinary addition discards every coefficient
of the left operand below the leading exponent of the right operand before
combining the leading coefficients. Implementing both additions separately
is essential to the tests.

\subsection{Recorded verification}
The completed checks are summarized below. Their precise grids, seed, and
sampling list are recorded in the JSON file.

\begin{center}
\renewcommand{\arraystretch}{1.22}
\begin{tabular}{@{}p{0.52\textwidth}r p{0.22\textwidth}@{}}
\toprule
Check & Cases & Result\\
\midrule
Finite Laurent formula versus sign expansion & 117,649 & All equal\\
Finite Laurent product inequality, ordered pairs & 15,625 & All passed\\
Strictness with a negative tail, within those pairs & 12,200 & All strict\\
Random dyadic-exponent product inequality & 100,000 & All passed\\
Dyadic real length versus actual sign length & 405 & All equal\\
Named birthday regressions & 12 & All passed\\
Ordinal arithmetic regressions & 3 & All passed\\
\bottomrule
\end{tabular}
\end{center}

The formula test uses exponents $-3,-2,-1,0,1,2$ and coefficient choices
\[
 \{0,1,-1,\tfrac12,2,\tfrac13,-\tfrac13\},
\]
so it tests all $7^6=117{,}649$ patterns, including zero coefficients and
missing-predecessor cases. The product test uses all $5^3=125$ polynomials on
exponents $-1,0,1$ with coefficients in
$\{0,1,-1,\tfrac12,\tfrac13\}$, and tests all ordered pairs. Exactly 881 pairs
attain equality; the strictness test concerns the specified 12,200-pair subset.

The exploratory search uses exponents $i/8$ for $-24\leq i\leq24$, one through
seven distinct exponents per factor, and rational coefficients drawn from the
ordered sampling list generated by
\[
 i/j,\qquad j\in(1,2,4,3,5),\quad i\in(-3,-2,-1,1,2,3).
\]
Repeated rational values in this list are intentionally retained. The random
seed is 394230. Among its 100,000 pairs, 202 attain equality. This search is
broader than the theorem in its allowed finite exponents, but much narrower
in its coefficients and support sizes.

\subsection{Running the checks}
From the extracted archive directory, run
\begin{verbatim}
python verify.py --output verification_results.json
\end{verbatim}
A smaller smoke test is available as
\begin{verbatim}
python verify.py --quick --output quick_results.json
\end{verbatim}
The output records the Python version and run date. Wall-clock duration is
machine-dependent. Reproduction should compare the mathematical counts and
outcomes, not timestamps or timing fields.

The program is not a formal verification of the proofs. It cannot quantify
over arbitrary real coefficients, infinite series, or all ordinals. The
argument for those cases is the mathematical one in the article, with its
external inputs explicitly identified.

\section{Proof audit and conclusion}
\label{sec:audit}
The proof has two nontrivial interfaces to standard surreal theory: the
normal-form/Hahn-series correspondence, and the sign-expansion theorem with
its two deletion rules. The additive birthday bound was recalled with its
usual recursive proof. No instance of the conjectured general multiplicative
bound is assumed in proving the Laurent case.

The internal chain of implications is
\[
\begin{gathered}
 \text{sign expansions}\ \Longrightarrow\
 \text{exact finite Laurent birthdays}\\
 \Longrightarrow\ \text{bounded and cross-product estimates}\\
 \Longrightarrow\ \text{finite Laurent product bound}\\
 \Longrightarrow\ \text{fixed-bound finite-jet argument}\\
 \Longrightarrow\ \text{full Laurent-series product bound}.
\end{gathered}
\]

Several potential gaps were checked explicitly. Negative coefficients reverse
signs but not lengths. An absent coefficient is dyadic zero, so it must not
activate the boundary correction. Adjacent nondyadic coefficients require the
second deletion rule even though the shared-prefix rule has already been
applied. Multiplying by a dyadic scalar can turn a nondyadic coefficient into
a dyadic one, but cannot turn a dyadic coefficient into a nondyadic one.
The only sharp cross-product case has a coefficient of birthday one and a
bounded tail with leading birthday coefficient one; it was treated separately.
Finally, infinite support was handled by a uniform bound on finite product
truncations, not by an unproved continuity property of natural operations.

The resulting affirmative restriction is the entire field
$\R((\omega^{-1}))$, not merely a collection of monomials or rational functions.
The exact finite and infinite birthday formulas supply additional information
that is independent of the product inequality. In particular they expose why
some straightforward-looking normal-form arguments lose the required bound,
and they identify a concrete next obstacle: the shared-prefix combinatorics of
fractional dyadic exponents. This is the extent of the result established in
this article; the full birthday product conjecture remains beyond this proof.

\appendix
\section{Bibliographic and provenance notes}
The research status discussion is deliberately narrower than a claim that a
complete literature survey has proved absence of later work. The 2001 source
contains the monomial theorem and the general weaker estimate; the 2024
statement is first-person testimony by one of its authors. The 2011 paper and
2022 preprint provide accessible formulations of the sign-expansion rules used
here. No purported solution of a different conjecture bearing Gonshor's name
was treated as a resolution of the birthday question.

The original paper by van den Dries and Ehrlich has an erratum
\cite{vdDEErratum}. The present proof does not invoke their estimates for
additive monoids of supports, closure under inversion, or exponential-field
constructions. The affirmative monomial theorem is cited for context; our
restricted product proof instead uses the elementary real-coefficient lemma
and the stated sign rules. This separation keeps the logical dependencies of
the argument visible.

The archive contains the article source and PDF, exact verification code,
recorded finite checks, a README with build instructions, and an audit file.
It does not include third-party books or papers, font files, or a claim of
proof-assistant certification.

\clearpage
\begin{thebibliography}{9}
\bibitem{Gonshor1986}
H.~Gonshor,
\emph{An Introduction to the Theory of Surreal Numbers},
London Mathematical Society Lecture Note Series, vol.~110,
Cambridge University Press, 1986.
Sign expansions: Chapter~5; birthday product conjecture: p.~96.

\bibitem{vdDE2001}
L.~van den Dries and P.~Ehrlich,
Fields of surreal numbers and exponentiation,
\emph{Fundamenta Mathematicae} \textbf{167} (2001), no.~2, 173--188.
\href{https://doi.org/10.4064/fm167-2-3}{doi:10.4064/fm167-2-3}.

\bibitem{vdDEErratum}
L.~van den Dries and P.~Ehrlich,
Erratum to ``Fields of surreal numbers and exponentiation,''
\emph{Fundamenta Mathematicae} \textbf{168} (2001), no.~3, 295--297.

\bibitem{Ehrlich2011}
P.~Ehrlich,
Conway names, the simplicity hierarchy and the surreal number tree,
\emph{Journal of Logic \& Analysis} \textbf{3} (2011), article~1, 1--26.
\href{https://doi.org/10.4115/jla.2011.3.1}{doi:10.4115/jla.2011.3.1}.
In particular, Propositions~3.6--3.8.

\bibitem{BG2022}
O.~Bournez and Q.~Guilmant,
Surreal fields stable under exponential and logarithmic functions,
arXiv:2201.08199v1, 20 January 2022.
\url{https://arxiv.org/abs/2201.08199v1}.
In particular, Definition~2.6 and Theorem~2.7.

\bibitem{Ehrlich2024}
P.~Ehrlich,
Answer to ``Birthday of combinatorial game product,''
MathOverflow, 13 August 2024.
\url{https://mathoverflow.net/questions/476824/}.
Accessed 20 September 2026.
\end{thebibliography}
\end{document}
